% ============================================================
% CAPITOLUL 10 — CONCLUZII ȘI DIRECȚII VIITOARE
% ============================================================

\chapter{Concluzii și Direcții de Dezvoltare Viitoare}
\label{ch:concluzii}

\section{Sinteza lucrării}
\label{sec:sinteza_contributii}

Lucrarea a fost reorientată în jurul modelului SIDARTHE, deoarece
acesta descrie mai bine situațiile în care diagnosticul, simptomatologia,
severitatea clinică și mortalitatea influențează evoluția observată a unei
epidemii. Modelele SIR/SIS/SEIR/SIRS rămân prezente, dar rolul lor este
secundar: oferă repere matematice și numerice simple, utile pentru comparații și
verificări de regresie.

Pe această bază a fost construit un sistem software în care datele, simulările și
metricile de performanță sunt tratate unitar. API-ul Go oferă autentificare,
organizare pe studii, control RBAC, încărcare de seturi de date și pornire de
simulări; Redis decuplează cererile utilizatorului de execuția numerică;
runner-ele R, Julia și Python rezolvă modelul cerut prin
\texttt{model\_type}; cele două interfețe Vue.js separă fluxurile de lucru ale
cercetătorilor de operațiile globale ale super-adminilor. Contribuția practică a
sistemului constă în această legătură între model, execuție, izolare
organizațională și observabilitatea rezultatelor.

Lucrarea a fost extinsă cu un cadru de comparație numerică și calibrare. Această
extensie clarifică modul în care rezultatele modelelor pot fi raportate la date
observate prin operatori de observație, metrici de eroare și validare pe ferestre
temporale separate. În plus, a fost introdusă o direcție predictivă bazată pe
LSTM, tratată ca extensie statistică pentru prognoză pe termen scurt și pentru
analiza reziduurilor rămase după simularea compartimentală.

\section{Concluzii privind modelele}
\label{sec:concluzii_modele}

SIDARTHE introduce o structură mai apropiată de realitatea raportării
epidemiologice decât SIR. Separarea cazurilor nedepistate de cele diagnosticate,
a formelor ușoare de cele severe și a deceselor permite formularea
unor întrebări care nu pot fi abordate natural într-un model cu trei
compartimente. În același timp, această granularitate crește numărul de parametri
și face calibrarea mai dificilă.

Modelele SIR, SIS, SEIR și SIRS rămân utile tocmai prin simplitate. Ele
permit verificarea conservării populației, a pozitivității soluției și a
comportamentului solverelor într-un cadru ușor de controlat. În sistemul
implementat, aceste modele sunt tratate ca repere secundare, nu ca explicații
complete ale datelor reale.

Comparația numerică trebuie interpretată împreună cu structura modelului. Un
RMSE mai mic nu este suficient pentru a declara un model superior dacă acesta
necesită mai mulți parametri neobservabili sau dacă supraînvață fereastra de
calibrare. Pentru SIDARTHE, valoarea reală este dată de capacitatea de a formula
întrebări clinice mai bogate, nu doar de reducerea unei erori agregate.

\section{Concluzii privind implementarea}
\label{sec:concluzii_implementare}

Implementarea arată că același contract de simulare poate fi rulat în ecosisteme
diferite fără a schimba API-ul sau interfața web. R oferă o formulare compactă
prin \texttt{deSolve}; Python oferă integrare directă cu infrastructura și
bibliotecile științifice uzuale; Julia este potrivită pentru experimente numerice
mai ambițioase, mai ales atunci când costul de compilare este amortizat.

Schimbarea schemei rezultatelor din coloane fixe \texttt{s/i/r} într-un obiect
JSONB \texttt{compartments} este una dintre modificările importante ale pivotării.
Ea face sistemul mai flexibil și permite adăugarea unor modele noi fără migrații
ale tabelei de rezultate. Costul acestei flexibilități este o serializare mai
generală și o nevoie mai mare de validare la nivel de aplicație.

Pe latura operațională, sistemul a fost deployat efectiv pe un cluster Kubernetes
v1.31 găzduit pe o platformă cloud privată, cu trei noduri control plane și trei
workeri, YugabyteDB la factor de replicare~3 și Helm pentru ciclul de viață al
bazei de date. Această experiență confirmă portabilitatea arhitecturii și pune
în evidență câteva ajustări specifice mediilor distribuite, precum aplicarea
migrațiilor prin \texttt{psql} în locul \texttt{golang-migrate} sau adăugarea
extensiei \texttt{pgcrypto} la momentul instalării, datorită bazei PostgreSQL~11
a YugabyteDB. Aceste ajustări sunt locale și nu afectează nici codul aplicativ,
nici interfața REST.

\section{Limitări}
\label{sec:limitari}

Prima limitare este calibrarea parametrilor SIDARTHE. Datele OWID folosite în
lucrare nu oferă direct distribuția pe compartimentele $D$, $A$, $R$, $T$ și $E$.
Inițializarea folosită de sistem este reproductibilă, dar nu înlocuiește o
estimare epidemiologică dedicată. Pentru rezultate interpretate clinic, modelul
trebuie calibrat pe date despre testare, spitalizare, terapie intensivă și
vindecări.

Calibrarea introdusă în lucrare este formulată metodologic, dar nu este încă un
flux complet automatizat în aplicație. Extinderea practică ar necesita un tip
nou de job, stocarea parametrilor calibrați și raportarea separată a erorii de
antrenare și de test.

A doua limitare privește evaluarea empirică. Infrastructura de benchmark este
implementată, dar concluziile cantitative necesită rulări repetate pe configurații
hardware documentate. Fără aceste campanii complete, analiza performanței rămâne un
cadru de măsurare și nu o demonstrație statistică a superiorității unui limbaj.

Componenta LSTM este implementată operațional printr-un runner dedicat, dar nu
este încă validată experimental ca rezultat predictiv final. Ea necesită serii
temporale suficiente, separare strict cronologică între antrenare și test și
comparație cu linii de bază simple, precum persistența sau media mobilă.

O altă limitare este nivelul de maturitate al aplicației web. Fluxurile principale
sunt acoperite, inclusiv studiile, rolurile și consola super-admin, însă lipsesc
încă teste automate end-to-end, mecanisme de retry mai sofisticate pentru job-uri,
politici de rate limiting și audit al operațiilor utilizatorilor. Aceste elemente
ar fi necesare înainte de folosirea sistemului într-un context operațional.

\section{Direcții viitoare}
\label{sec:directii_viitoare}

Direcția imediată este calibrarea modelului SIDARTHE pe date mai bogate decât
seria de cazuri active, în special date despre testare, spitalizare, terapie
intensivă, vindecări și decese. În plus, modelele SIR/SIS/SEIR/SIRS pot fi
folosite pentru scenarii de bază și pentru testarea
sensibilității la ipoteze epidemiologice mai simple.

O direcție tehnică imediată este implementarea calibrării ca job asincron în
același mecanism Redis folosit pentru simulări. Rezultatul ar trebui să includă
parametrii estimați, erorile MAE/RMSE pe train și test, numărul de evaluări ale
solverului și timpul total de optimizare.

Pentru componenta ML, următorul pas realist este evaluarea runner-ului LSTM
dedicat împotriva unei linii de bază de persistență și a modelului
compartimental calibrat. Rezultatul ar trebui raportat prin MAE, RMSE, nRMSE,
timp de antrenare și stabilitate pe fereastra de test.

La nivel de sistem, următorul pas este completarea campaniilor de benchmark și
adăugarea unei suite automate de teste. Pentru backend, aceasta înseamnă teste
unitare și de integrare cu o bază de date izolată; pentru frontend, teste care
verifică fluxurile reale din browser. În zona de infrastructură, ar fi utilă
introducerea unor mecanisme explicite de retry și dead-letter queue pentru job-urile
care eșuează repetat, precum și activarea autoscalării orizontale (HPA) a
runner-elor pe metrici personalizate (de exemplu lungimea cozii Redis) și a
politicilor de rețea (NetworkPolicy) pentru izolare zero-trust intra-cluster.
Terminarea TLS automată prin \texttt{cert-manager} și un stack de
observabilitate (Prometheus + Grafana) ar completa profilul operațional al
deployment-ului Kubernetes.

Pe termen mai lung, sistemul poate fi extins cu vizualizări geografice, surse de
date actualizate periodic și module de alertare. Aceste extensii trebuie construite
pe aceeași disciplină metodologică: ipotezele modelului, calitatea datelor și
metricile software trebuie să rămână vizibile.
